\documentclass[10pt, aspectratio=169]{beamer}

\usepackage{clases}

\title{Investigación Operativa - Clase 4}
\subtitle{Minimax y maximin}
\date{}
\author{Nazareno Faillace Mullen}
\institute{Departamento de Matemática - FCEN - UBA}

\begin{document}

\maketitle


\begin{frame}{Para ir pensando}
    ¿Cómo modelarías linealmente el siguiente problema?

    \[
        \begin{array}{rlcl}
            \min & \multicolumn{3}{l}{\max\{2x_1+x_2, -x_1+3x_2\}}  \\
            \text{s.a:} & x_1 + x_2 & \leq & 4 \\
             &          2x_1-x_2 & \geq & 2 \\
             &          x_1, x_2 &\in & \mathbb{R}_{\geq 0}
        \end{array}
    \]
\end{frame}

\begin{frame}{Minimax}
    Lo reformulamos como un problema de programación lineal añadiendo una nueva variable real $w$:
    \[
        \begin{array}{rlcl}
            \min & \multicolumn{3}{l}{\color{red}w}  \\
            \text{s.a:} & x_1 + x_2 & \leq & 4 \\
             &          2x_1-x_2 & \geq & 2 \\
             &          \color{red}2x_1+x_2 & \color{red}\leq & \color{red}w \\
             &          \color{red}-x_1+3x_2 & \color{red}\leq & \color{red}w \\
             &          x_1, x_2 &\in & \mathbb{R}_{\geq 0} \\
             &          \color{red}w \in \color{red}\mathbb{R}
        \end{array}
    \]
    La solución óptima es $(x_1^\ast, x_2^\ast, w^\ast) = (1, 0, 2)$
\end{frame}

\begin{frame}{Minimax}
    \begin{itemize}
        \item El objetivo es minimizar el máximo de un conjunto de funciones lineales, sujetos a ciertas restricciones
    \end{itemize}
    \[\begin{array}{crl}
        \min & \multicolumn{2}{l}{\max\{f_1(x), f_2(x),\dots,f_\ell(x)\}}\\
        s.a.: & Ax\leq b \\
%         & x\geq0
    \end{array}
    \]
\end{frame}

\begin{frame}{Minimax - transformación a problema lineal}
    \begin{itemize}
        \item Consideramos una nueva variable real $w$ y le imponemos el siguiente conjunto de restricciones:
        \[w \geq f_k(x) \quad \forall k\in\{1,\dots,\ell\} \]
        El problema queda formulado entonces como:
        \[\begin{array}{crl}
        \min & \multicolumn{2}{l}{w}\\
        s.a.: & Ax &\leq b \\
         & w &\geq f_k(x) \quad \forall k\in\{1,\dots,\ell\} \\
%         & x &\geq0 \\
         & x &\in \mathbb{K}^n\\
         & w &\in \mathbb{R}
        \end{array}
        \]
    \end{itemize}
\end{frame}

\begin{frame}{Maximin}
    \textbf{Obs:} una idea análoga sirve para el caso maximin

    \begin{itemize}
        \item El objetivo es maximizar el mínimo de un conjunto de funciones lineales, sujetos a ciertas restricciones
    \end{itemize}
    \[\begin{array}{crl}
        \max & \multicolumn{2}{l}{\min\{f_1(x), f_2(x),\dots,f_\ell(x)\}}\\
        s.a.: & Ax\leq b \\
%         & x\geq0
    \end{array}
    \]
\end{frame}


\begin{frame}{Maximin - transformación a problema lineal}
    \begin{itemize}
        \item Consideramos una nueva variable real $w$ y le imponemos el siguiente conjunto de restricciones:
        \[w \leq f_k(x) \quad \forall k\in\{1,\dots,\ell\} \]
        El problema queda formulado entonces como:
        \[\begin{array}{crl}
        \max & \multicolumn{2}{l}{w}\\
        s.a.: & Ax &\leq b \\
         & w &\leq f_k(x) \quad \forall k\in\{1,\dots,\ell\} \\
%         & x &\geq0 \\
         & x &\in \mathbb{K}^n\\
         & w &\in \mathbb{R}
        \end{array}
        \]
    \end{itemize}
\end{frame}

\begin{frame}{Ejercicio}
    \scriptsize
    \parbox{\textwidth}{Para dar a conocer sus productos, MicroApple decidió invertir  $850$
        mil pesos en publicidad. Hay tres proyecciones distintas de audiencia (en miles de personas)
        por cada anuncio en cada uno de los medios:
    \begin{table}[]
        \centering
        \begin{tabular}{ccccc}\toprule
             & Redes Sociales & TV & Radio & Periódico \\ \midrule
            P1 & 50  & 31 & 20 & 26 \\
            P2 & 32  & 41 & 15 & 19 \\
            P3 & 60  & 38 & 22 & 25
        \end{tabular}
    \end{table}

        Los datos se encuentran en la siguiente tabla:
        \begin{table}[]
            \centering
            \begin{tabular}{ccccc}
                \toprule
                & Redes Sociales & TV & Radio & Periódico \\ \midrule
                %Audiencia (en miles) & 50 & 31 & 20 & 35 \\
                Costo por anuncio (en miles) & 35 & 29    & 22     & 27        \\
                Máxima cantidad de anuncios  & 20 & 15    & 25     & 15
            \end{tabular}
        \end{table}
        Además de respetar el presupuesto, el departamento de marketing requiere que:
        \begin{enumerate}[(1)]
            \item
            el valor absoluto de la diferencia entre la cantidad de anuncios en medios masivos (redes soluciones y
            televisión) y la cantidad en medios tradicionales (radio y periódico) sea a lo sumo $10$.
            \item los anuncios en radio deben ser al menos el $5$\% de los anuncios totales
        \end{enumerate}

        Elaborar un modelo para decidir cuántos anuncios de cada medio debe comprar MicroApple para maximizar el mínimo de
        audiencia alcanzada según las proyecciones.
    }
\end{frame}


\end{document}



